\providecommand{\Obs}{\mathrm{Obs}}
\providecommand{\hCS}{\mathrm{hCS}}
\providecommand{\Ch}{\mathrm{Ch}}
\providecommand{\Dolb}{\mathrm{Dolb}}
\providecommand{\Conf}{\mathrm{Conf}}
\providecommand{\Fact}{\mathrm{Fact}}
\providecommand{\HH}{\mathrm{HH}}
\providecommand{\Def}{\mathrm{Def}}
\providecommand{\CE}{\mathrm{CE}}
\providecommand{\Sym}{\mathrm{Sym}}
\providecommand{\Hol}{\mathrm{Hol}}
\providecommand{\Tr}{\operatorname{Tr}}
\providecommand{\Rep}{\mathrm{Rep}}
\providecommand{\Alg}{\mathrm{Alg}}
\providecommand{\MC}{\mathrm{MC}}
\providecommand{\End}{\operatorname{End}}
\providecommand{\BV}{\mathrm{BV}}
\providecommand{\dbar}{\bar\partial}
\providecommand{\C}{\mathbb{C}}
\providecommand{\cO}{\mathcal{O}}
\providecommand{\cE}{\mathcal{E}}
\providecommand{\cA}{\mathcal{A}}
\providecommand{\fg}{\mathfrak{g}}
\providecommand{\fh}{\mathfrak{h}}
\providecommand{\fa}{\mathfrak{a}}
\providecommand{\Eone}{E_1}
\providecommand{\Etwo}{E_2}
\providecommand{\Ethree}{E_3}
\providecommand{\PhiFA}{\Phi^{\mathrm{FA}}}
\providecommand{\SpCh}{\mathrm{Sp}^{\mathrm{ch}}}
\providecommand{\Zchder}{Z^{\mathrm{der}}_{\mathrm{ch}}}

\section*{Deformation Complex of Six-Dimensional hCS on $\C^3$}

\paragraph{The local hCS datum.}
Let \(X=\C^3\) with holomorphic volume form
\[
  \Omega=dz_1\wedge dz_2\wedge dz_3,
\]
and let \((\fg,\langle-,-\rangle)\) be a finite-dimensional quadratic
complex Lie algebra. On a relatively compact open \(U\subset X\), the
classical holomorphic Chern--Simons local Lie algebra is
\[
  \fh_{\hCS}(U)
  =
  \Omega^{0,\bullet}_c(U,\fg)[1],
  \qquad
  \ell_1=\dbar,
  \qquad
  \ell_2(\alpha,\beta)=[\alpha\wedge\beta]_{\fg},
  \qquad
  \ell_n=0\quad(n\geq3).
\]
The Maurer--Cartan equation is
\[
  \dbar A+\frac12[A,A]=0,
\]
the integrability equation for the \(\dbar\)-operator
\(\dbar+A\). The classical observables form the Chevalley--Eilenberg
cochain factorisation algebra
\[
  \Obs^{\mathrm{cl}}_{\hCS}(U)
  =
  C^\bullet_{\CE,\mathrm{cont}}\bigl(\fh_{\hCS}(U)\bigr)
  =
  \widehat{\Sym}\bigl(\fh_{\hCS}(U)^\vee[-1]\bigr),
  \qquad
  d_{\Obs}^{\mathrm{cl}}=d_{\dbar}+d_{\ell_2}.
\]

\paragraph{Quantum observables and the holomorphic \(\Ethree\)-level.}
A quantum hCS factorisation algebra is the preceding classical
factorisation algebra equipped with a Costello renormalisation datum:
an elliptic gauge fixing, heat kernels, regularised propagators,
counterterms, completed local functionals, and effective interactions
\(I[L]\) satisfying
\[
  QI[L]+\frac12\{I[L],I[L]\}_L+\hbar\Delta_LI[L]=0.
\]
When the quantum master equation is solved, the scale-\(L\) observables
are
\[
  \Obs^q_{\hCS,L}(U)
  =
  \left(
  \Obs^{\mathrm{cl}}_{\hCS}(U)[[\hbar]],
  d_{\Obs}^{\mathrm{cl}}+\{I[L],-\}_L+\hbar\Delta_L
  \right).
\]
The factorisation operations over configurations in \(\C^3\) give a
holomorphic \(\Ethree\)-factorisation algebra in \(\Ch(\Dolb)\). A
translation-invariant Dolbeault contraction supplies an ordinary
\(\Ethree\)-algebra model. The contraction is extra data; the
holomorphic factorisation object is the native local object.

\paragraph{Formal deformation complex.}
Let \(A\) be either the solved quantum observable algebra
\(\Obs^q_{\hCS}\), or the classical observable algebra when the
quantum datum has not been chosen. In a stable symmetric monoidal
ambient category, the formal deformation complex of \(A\) as an
\(\Ethree\)-algebra is
\[
  \Def_{\Ethree}(A)
  =
  \HH^\bullet_{\Ethree}(A,A)[3].
\]
The shift is the Francis--Lurie \(E_n\)-tangent shift at \(n=3\).
This complex carries a canonical \(L_\infty\)-structure. Its
cohomology has the following deformation-theoretic meaning:
\[
  H^1\Def_{\Ethree}(A)
  \quad\hbox{classifies first-order deformations,}
  \qquad
  H^2\Def_{\Ethree}(A)
  \quad\hbox{contains obstruction classes.}
\]
A first-order cocycle \(\alpha_1\) extends to a formal deformation
\(\alpha(t)=t\alpha_1+t^2\alpha_2+\cdots\) only when the Maurer--Cartan
equation
\[
  d_{\Def}\alpha
  +\frac12\ell_2^{\Def}(\alpha,\alpha)
  +\frac1{3!}\ell_3^{\Def}(\alpha,\alpha,\alpha)
  +\cdots
  =0
\]
can be solved recursively. The primary second-order obstruction is the
class of \(\frac12\ell_2^{\Def}(\alpha_1,\alpha_1)\) in
\(H^2\Def_{\Ethree}(A)\); higher brackets give the subsequent
obstructions. Thus an explicit element of
\(\HH^\bullet_{\Ethree}(A,A)[3]\) is a first-order deformation class
until these obstruction equations are checked.

\paragraph{Abelian affine computation.}
For abelian \(\fg\), the bracket \(\ell_2\) vanishes. The hCS action is
quadratic,
\[
  S_2(\cA)
  =
  \frac12\int_{\C^3}\Omega\wedge
  \langle\cA,\dbar\cA\rangle,
\]
and the Costello quantum master equation has no interaction-graph
anomaly. With polynomial boundary conditions the Dolbeault cohomology is
\[
  H^\bullet_{\dbar}(\C^3,\fg)
  =
  \C[z_1,z_2,z_3]\otimes\fg
  \quad\hbox{in Dolbeault degree }0.
\]
With compact support the one-dimensional Serre-dual class is
\[
  H^{0,3}_{\dbar,c}(\C^3)\cong\C.
\]
After choosing the standard affine Dolbeault contraction, the
translation-invariant polynomial first-order sector of
\(\Def_{\Ethree}(\Obs^q_{\hCS})\) is represented by
\[
  \left(
  \Omega^{0,\bullet}(\C^3)
  \otimes
  \widehat{\Sym}\bigl(\fg^\vee[-1]\oplus\fg[-2]\bigr),
  \dbar
  \right),
\]
with the displayed shifts coming from the BV field-antifield pairing
and the \(E_3\)-Hochschild shift. This model covers the controlled
abelian affine sector. Compact and nonabelian hCS deformations require
the additional obstruction theory stated below.

\paragraph{The compactly supported volume class.}
Choose
\(\psi\in\Omega^{0,3}_c(\C^3)\) with
\(\int_{\C^3}\Omega\wedge\psi=1\). The class \([\psi]\) gives a
degree-one local cochain in the abelian deformation complex by inserting
the pointwise BV contraction against the density \(\Omega\wedge\psi\):
\[
  \alpha_\psi(F)
  =
  \int_{\C^3}\Omega\wedge\psi(z)\,
  \Delta^{\mathrm{vert}}_zF .
\]
Here \(\Delta^{\mathrm{vert}}_z\) is the regularised second-variation
operator contracting the field and antifield variables at \(z\). Since
\(\dbar\psi=0\) and the free heat-kernel BV operator commutes with
\(\dbar\), \(d_{\Def}\alpha_\psi=0\). The first obstruction is
\[
  \left[
  \frac12\ell_2^{\Def}(\alpha_\psi,\alpha_\psi)
  \right]
  \in H^2\Def_{\Ethree}(\Obs^q_{\hCS}).
\]
In the constant-coefficient abelian subcomplex, the vertical
contractions commute after regularisation, so this obstruction class
vanishes and the one-parameter differential
\[
  d_{\Obs^q}+t\alpha_\psi
\]
defines a formal deformation inside that subcomplex. This proves
unobstructedness for the displayed scalar abelian family. Arbitrary hCS
deformations, nonabelian interactions, and compact Calabi--Yau
threefolds remain governed by their own obstruction brackets.

\paragraph{Nonabelian anomaly slot.}
Let \(\rho:\fg\to\End(V)\) be a finite-dimensional representation. The
dimension-three gauge anomaly is governed by the symmetric quartic
invariant
\[
  I_4^\rho(x_1,x_2,x_3,x_4)
  =
  \frac1{4!}
  \sum_{\sigma\in S_4}
  \Tr_V\!\left(
  \rho(x_{\sigma(1)})\rho(x_{\sigma(2)})
  \rho(x_{\sigma(3)})\rho(x_{\sigma(4)})
  \right)
  \in\Sym^4(\fg^\vee)^\fg .
\]
The corresponding local Lie algebra cocycle has representative
\[
  \Theta^\rho_U(\alpha_0,\alpha_1,\alpha_2,\alpha_3)
  =
  \int_U
  \left[
  I_4^\rho\bigl(
  \alpha_0,\partial\alpha_1,\partial\alpha_2,\partial\alpha_3
  \bigr)
  \right]_{(3,3)},
  \qquad
  \alpha_i\in\Omega^{0,\bullet}_c(U,\fg).
\]
Its cohomology class is the degree-one obstruction to solving the
quantum master equation.

The quadratic Casimir \(C_2^\rho\) belongs to the two-point
field-strength counterterm,
\[
  S_{\mathrm{ct}}^{(2)}
  =
  -\hbar\,\frac{C_2^\rho}{(4\pi)^3}
  \log(L/\varepsilon)
  \int_U\Omega\wedge\Tr_\rho(\cA\,\dbar\cA),
\]
and is separate from the quartic quantum-master-equation obstruction.

\paragraph{Compact closure hypotheses.}
For a compact Calabi--Yau threefold \(Y\), a nonabelian quantum hCS
factorisation algebra requires more than the local first-order
deformation complex. The data are:
\[
  [\Theta^\rho_Y]=0
  \quad\hbox{or a counterterm }S_{\mathrm{ct}}\hbox{ with }
  Q_{\BV}S_{\mathrm{ct}}=-\Theta^\rho_Y,
\]
an orientation and determinant-line framing, a holomorphic volume form,
elliptic gauge fixing, heat-kernel renormalisation, RG-compatible
counterterms, factorisation or TCFT sewing, and analytic completion of
the OPE operations. Under these hypotheses the same
\(\Def_{\Ethree}=\HH^\bullet_{\Ethree}(-,-)[3]\) controls formal
deformations of the solved observable algebra. Without these
hypotheses the local complex supplies candidate first-order directions
and obstruction groups.

\paragraph{Specialisation to the CY-to-chiral output.}
The bulk hCS factorisation object belongs to the first stage of the
CY-to-chiral construction:
\[
  \PhiFA_3(\mathcal C)\in
  \Ethree\hbox{-}\mathrm{HolFA}(X).
\]
For an admissible specialisation pair \((\Sigma_2,C)\), the
dimension-three chiral output is
\[
  A_C
  =
  \SpCh_{\Sigma_2,C}\bigl(\PhiFA_3(\mathcal C)\bigr)
  \in
  \Alg^{\mathrm{ch}}_{\Eone}(C).
\]
The braided \(\Etwo\)-layer is carried by centre, double, or module
constructions, for example
\[
  Z\bigl(\Rep^{\Eone}(A_C)\bigr)
  \simeq
  \Rep^{\Etwo}\bigl(
    \Zchder(A_C)
  \bigr)
\]
on the constructed Drinfeld-centre locus. Deforming the bulk
\(\Ethree\)-factorisation algebra and deforming the specialised
\(\Eone\)-chiral algebra are connected by the specialisation functor,
with separate obstruction complexes on the two sides.

\paragraph{Chain-level and \((\infty,1)\)-categorical forms.}
The chain-level statement is the explicit Dolbeault/BV complex:
\[
  \Omega^{0,\bullet}(\C^3,\fg)[1],
  \qquad
  d=\dbar,
  \qquad
  \omega_{\BV}(\alpha,\beta)
  =
  \int_{\C^3}\Omega\wedge\langle\alpha,\beta\rangle,
\]
together with the heat-kernel construction of \(\Delta_L\), the
cochain \(\alpha_\psi\), and the obstruction brackets
\(\ell_m^{\Def}\). The \((\infty,1)\)-categorical statement is that the
derived moduli problem of \(\Ethree\)-algebras in \(\Ch(\Dolb)\) has
tangent \(L_\infty\)-algebra
\(\HH^\bullet_{\Ethree}(A,A)[3]\) at \(A\). These two statements have
equal force: the first supplies representatives and computations; the
second supplies functoriality, base-change, and obstruction-theoretic
meaning.

\paragraph{References.}
\begin{itemize}
  \item K.~Costello, \emph{Renormalization and Effective Field Theory}
  (2011), Chapter~2 and Section~5.9, for heat-kernel regularisation,
  the BV Laplacian, and quantum-master-equation obstruction theory.
  \item K.~Costello and O.~Gwilliam, \emph{Factorization Algebras in
  Quantum Field Theory}, Vols.~I--II, for local dg Lie algebras,
  classical and quantum observables, and factorisation descent.
  \item K.~Costello and S.~Li, holomorphic BV anomaly analysis, for the
  quartic invariant \(I_4^\rho\) in complex dimension three.
  \item J.~Francis, ``The tangent complex and Hochschild cohomology of
  \(E_n\)-rings,'' \emph{Compos.\ Math.} 149 (2013), for the
  \(E_n\)-Hochschild tangent complex and the \(n\)-shift.
  \item J.~Lurie, \emph{Higher Algebra}, Section~5.3, for formal
  moduli of \(E_n\)-algebras in stable symmetric monoidal categories.
  \item O.~Gwilliam and B.~Williams, ``Higher Kac--Moody algebras and
  symmetries of holomorphic field theories,'' for holomorphic
  factorisation algebras and holomorphic \(E_n\)-structures.
\end{itemize}
